\documentclass[11pt,a4paper]{article}
\usepackage[margin=1in]{geometry}
\usepackage{amsmath,amssymb,amsthm}
\usepackage{mathtools}
\usepackage{booktabs}
\usepackage{enumitem}

\newtheorem{theorem}{Theorem}
\newtheorem{corollary}[theorem]{Corollary}
\newtheorem{proposition}[theorem]{Proposition}
\newtheorem{conjecture}[theorem]{Conjecture}
\theoremstyle{definition}
\newtheorem{definition}[theorem]{Definition}
\theoremstyle{remark}
\newtheorem{remark}[theorem]{Remark}

\newcommand{\CP}{\mathbb{CP}}
\newcommand{\C}{\mathbb{C}}
\newcommand{\R}{\mathbb{R}}
\newcommand{\Z}{\mathbb{Z}}
\newcommand{\SU}{\mathrm{SU}}
\newcommand{\su}{\mathfrak{su}}

\title{The Koide Canon\\[6pt]\large Complete Fermion Mass Structure from $H=3$}
\author{Extracted from the Irreducible Twistor Geometry Framework\\J.~R.~Manuel}
\date{April 2026}

\begin{document}
\maketitle

\section*{The Formula}

\begin{definition}[Koide mass parameterisation]
The charged fermion masses in each sector are
\begin{equation}\label{eq:koide}
\boxed{\;m_k \;=\; M^2 \left(1 + \sqrt{2}\,\cos\!\left(\theta + \frac{2\pi k}{3}\right)\right)^{\!2},\qquad k = 0,\,1,\,2\;}
\end{equation}
where $M^2$ is the sector mass scale, $\theta$ is the sector phase, and $k$ labels the three generations.
\end{definition}

Every parameter is determined by two numbers: $H=3$ (the unique positive root of $n(n\!-\!1)(n\!+\!2)(n\!+\!3)=0$) and $K^*=7/30$ (the equilibrium conflict from the conservation law).

\bigskip

\section*{The Parameters}

\renewcommand{\arraystretch}{1.4}
\begin{center}
\begin{tabular}{@{}llll@{}}
\toprule
\textbf{Parameter} & \textbf{Formula} & \textbf{Value} & \textbf{Status} \\
\midrule
$Q$ & $(H\!-\!1)/H$ & $2/3$ & Proved \\
$r$ & $\sqrt{6Q-2}$ & $\sqrt{2}$ & Proved (follows from $Q$) \\
$\theta_\ell$ (leptons) & $2/(nH^2),\; n\!=\!1$ & $2/9$ & Proved \\
$\theta_d$ (down quarks) & $2/(nH^2),\; n\!=\!H\!-\!1\!=\!2$ & $1/9$ & Proved \\
$\theta_u$ (up quarks) & $2/(nH^2),\; n\!=\!H\!=\!3$ & $2/27$ & Proved \\
$\theta_\nu$ (neutrinos) & $K^*/(H\!-\!1)$ & $7/60$ & Proved (dual formula) \\
\bottomrule
\end{tabular}
\end{center}

\bigskip

The Koide ratio across sectors (conjectured for quarks):
\begin{equation}\label{eq:Q_sectors}
Q \;=\; \frac{H\!-\!1}{H} \;+\; \frac{C}{H^2\!-\!1} \;+\; T_3\cdot\frac{K^*}{2}
\end{equation}
\begin{center}
\begin{tabular}{@{}lcccl@{}}
\toprule
\textbf{Sector} & $C$ & $T_3$ & $Q$ & \textbf{Status} \\
\midrule
Leptons & $0$ & $0$ & $2/3$ & Proved \\
Down quarks & $1$ & $-1/2$ & $11/15$ & Conjecture \\
Up quarks & $1$ & $+1/2$ & $17/20$ & Conjecture \\
\bottomrule
\end{tabular}
\end{center}

\section*{The Identities}

For any angle $\theta$, when $r=\sqrt{2}$:

\begin{align}
\sum_{k=0}^{2} \sqrt{x_k} &= 3 = H \label{eq:sum_sqrt}\\[4pt]
\sum_{k=0}^{2} x_k &= 6 = 2H \label{eq:sum_x}\\[4pt]
Q \;=\; \frac{\sum x_k}{\bigl(\sum\sqrt{x_k}\bigr)^2} &= \frac{6}{9} = \frac{2}{3} = \frac{H\!-\!1}{H} \label{eq:Q}
\end{align}

\emph{Proof.} The three phases $\phi_k = \theta + 2\pi k/3$ are equally spaced on the circle. By Fourier orthogonality on $\Z_3$: $\sum\cos\phi_k = 0$ and $\sum\cos^2\!\phi_k = 3/2$. Then $\sum\sqrt{x_k} = 3 + \sqrt{2}\cdot 0 = 3$, and $\sum x_k = 3 + 2\sqrt{2}\cdot 0 + 2\cdot 3/2 = 6$. \qed

\section*{The Dimensionless Mass Factors}

At $\theta = 2/9$ (leptons):
\begin{center}
\begin{tabular}{@{}lrl@{}}
\toprule
& $x_k = \bigl(1+\sqrt{2}\cos(2/9+2\pi k/3)\bigr)^2$ & \\
\midrule
$x_e$ & $0.001\,628\,12$ & (electron) \\
$x_\mu$ & $0.336\,645\,87$ & (muon) \\
$x_\tau$ & $5.661\,726\,01$ & (tau) \\
\midrule
$\sum x_k$ & $6.000\,000\,00$ & $= 2H$ (exact) \\
$(\sum\sqrt{x_k})^2$ & $9.000\,000\,00$ & $= H^2$ (exact) \\
\bottomrule
\end{tabular}
\end{center}

\section*{The Topological Cliff}

The mass factor $x_k$ vanishes when $1 + \sqrt{2}\cos\phi_k = 0$, i.e.\ at $\phi = 3\pi/4 = 135^\circ$. The electron sits at
\[
\phi_e = \frac{2}{9} + \frac{2\pi}{3} \approx 132.73^\circ.
\]
A margin of $2.27^\circ$ from the zero-mass cliff. The entire lepton hierarchy $m_\tau/m_e = 3477$ is the ratio of a constructive peak to this near-node, with the node's position locked by $\theta = 2/9$ from $H = 3$. The electron is light because the Coxeter number places it at the edge.

\section*{The Mass Scales}

\begin{center}
\begin{tabular}{@{}lccl@{}}
\toprule
\textbf{Sector} & $M^2/m(0^{++})$ & \textbf{Decomposition} & \textbf{Deviation} \\
\midrule
Leptons & $11/60$ & $\dfrac{1}{d_\mathrm{super}} + \dfrac{1}{H(H\!+\!1)} = \dfrac{1}{10}+\dfrac{1}{12}$ & $0.11\%$ \\[8pt]
Down quarks & $3/8$ & $\dfrac{1}{\dim(\SU(3))} + \dfrac{1}{H\!+\!1} = \dfrac{1}{8}+\dfrac{1}{4}$ & $1.3\%$ \\[8pt]
Up quarks & $40/3$ & Forced by product constraint & $0.1\%$ \\[4pt]
\bottomrule
\end{tabular}
\end{center}

\smallskip
Product constraint: $\;\dfrac{11}{60}\times\dfrac{3}{8}\times\dfrac{40}{3} = \dfrac{11}{12} = \dfrac{H(H\!+\!1)-1}{H(H\!+\!1)} = \text{massive/total generators}.$

\section*{The Quadratic}

The lepton denominators $10$ and $12$ are the roots of
\begin{equation}\label{eq:quadratic}
\boxed{\;x^2 - 22x + 120 = (x-10)(x-12) = 0\;}
\end{equation}
\begin{center}
\begin{tabular}{@{}ll@{}}
\toprule
\textbf{Coefficient} & \textbf{Meaning} \\
\midrule
Sum $= 22 = 2(H^2\!+\!H\!-\!1)$ & $= d_\mathrm{bos} - H - 1 = 26 - 4$ \\
Product $= 120 = H(H\!+\!1)(H^2\!+\!1) $ & $= S + 1/K^* = 5!$ \\
Discriminant $= 4 = (H\!-\!1)^2 = H\!+\!1$ & the self-consistency equation \\
\bottomrule
\end{tabular}
\end{center}
The Koide lepton scale is $\dfrac{11}{60} = \dfrac{\text{sum of roots}}{\text{product of roots}}$.

\medskip

The down quark denominators $8$ and $4$ satisfy $x^2 - 12x + 32 = 0$, nested: the down sum ($12$) is a lepton root, and the lepton discriminant ($4$) is a down root.

\section*{The Glueball--Lepton Sum Rule}

Since $\sum x_k = 6$ and $M^2 = m(0^{++})\times 11/60$:
\begin{equation}\label{eq:sum_rule}
\boxed{\;m(0^{++}) \;=\; \frac{10}{11}\bigl(m_e + m_\mu + m_\tau\bigr) \;=\; \frac{d_\mathrm{super}}{d_\mathrm{super}+1}\;\sum m_\ell\;}
\end{equation}
Numerically: $\frac{10}{11}\times 1883.0 = 1711.8\;\mathrm{MeV}$ (lattice QCD: $1710\pm 50$).

\begin{remark}[The beta function coincidence]
The numerator $11 = H(H\!+\!1)-1$ counts massive gauge bosons. The pure $\SU(3)$ one-loop beta coefficient is $\beta_0 = 11N_c/3 = 11$ at $N_c = H = 3$. Setting $H(H\!+\!1)-1 = 11H/3$ gives $3H^2 - 8H - 3 = 0$, with unique positive root $H = 3$. The running rate equals the particle content only at the Coxeter number.
\end{remark}

\section*{The Conversion System}

\emph{One measurement determines everything.}

\begin{enumerate}[nosep]
\item Pick any precisely known mass (e.g.\ $m_e$ or $m_\mu$, both known to $<0.001\%$).
\item Compute $M^2 = m_k/x_k$. From $m_e$: $M^2 = 313.859\;\mathrm{MeV}$. From $m_\mu$: $M^2 = 313.856\;\mathrm{MeV}$. Consistent to $0.001\%$.
\item Compute $m(0^{++}) = M^2\times 60/11 = 1712.0\pm 0.3\;\mathrm{MeV}$.
\item Every other mass follows: lepton masses from $x_k$, quark masses from the sector Koide formulae, hadron masses from the eigenvalue spectrum, electroweak masses from $H$-polynomial multipliers.
\end{enumerate}

\medskip

Zero dimensionless free parameters. One dimensional anchor (choice of energy unit).

\section*{Open Problems}

\begin{enumerate}[nosep]
\item \textbf{The $452\sigma$ residual.} The $e/\mu$ mass ratio deviates from $\theta=2/9$ by $9.83\times 10^{-6}$ ($452\sigma$ on the ratio, but only $0.97\sigma$ on $m_\tau$). Either the tau mass is ${\sim}1\sigma$ low, or the Koide formula carries a correction at $10^{-5}$. BESIII and Belle~II will decide.
\item \textbf{Quark $Q$ values.} The formula $Q = (H\!-\!1)/H + C/(H^2\!-\!1) + T_3\cdot K^*/2$ is conjectured, not derived, for colour-charged sectors.
\end{enumerate}

\end{document}
